%; whizzy chapter -dvi
% -initex iniptex -latex platex -format platex -bibtex jbibtex -fmt fmt
% 以上 whizzytex を使用する場合の設定。
 
%     Tokyo Debian Meeting resources
%     Copyright (C) 2012 Junichi Uekawa
%     Copyright (C) 2011, 2015 Nobuhiro Iwamatsu

%     This program is free software; you can redistribute it and/or modify
%     it under the terms of the GNU General Public License as published by
%     the Free Software Foundation; either version 2 of the License, or
%     (at your option) any later version.

%     This program is distributed in the hope that it will be useful,
%     but WITHOUT ANY WARRANTY; without even the implied warranty of
%     MERCHANTABILITY or FITNESS FOR A PARTICULAR PURPOSE.  See the
%     GNU General Public License for more details.

%     You should have received a copy of the GNU General Public License
%     along with this program; if not, write to the Free Software
%     Foundation, Inc., 51 Franklin St, Fifth Floor, Boston, MA  02110-1301 USA

%  preview (shell-command (concat "evince " (replace-regexp-in-string "tex$" "pdf"(buffer-file-name)) "&"))

%%ここからヘッダ開始。

\documentclass[mingoth,a4paper]{jsarticle}
\usepackage{monthlyreport}
% 日付を定義する、毎月変わります。
\newcommand{\debmtgyear}{2015}
\newcommand{\debmtgmonth}{11}
\newcommand{\debmtgdate}{21}
% started from zero:
% (let ((year 2013) (month 7)) (+ (* (- year 2005) 12) month -1))
\newcommand{\debmtgnumber}{134}

% tikz picture の為のマクロ設定
\usepackage[dvipdfmx]{graphicx}
\usepackage{tikz}

\begin{document}

\begin{titlepage}
\thispagestyle{empty}
% タイトルページ:編集必要な部分は最初のマクロに飛ばすこと

\vspace*{-2cm}
第\debmtgnumber{}回 東京エリア Debian 勉強会資料\\
\hspace*{-2cm}
\includegraphics{image2012-natsu/dotdeb.pdf}\\
\hfill{}\debmtgyear{}年\debmtgmonth{}月\debmtgdate{}日

% ここはアップデートすること
% 全角文字にしないとフォントのサイズが合わないので注意
\rotatebox{10}{\fontsize{30}{30} {\gt 特集：}}\\

\vspace*{-2cm}
\hfill{}\includegraphics[height=6cm]{image200502/openlogo-nd.eps}
\end{titlepage}

\newpage

\begin{minipage}[b]{0.2\hsize}
 \definecolor{titleback}{gray}{0.9}
 \colorbox{titleback}{\rotatebox{90}{\fontsize{80}{80} {\gt デビアン勉強会} }}
\end{minipage}
\begin{minipage}[b]{0.8\hsize}
\hrule
\vspace{2mm}
\hrule
\begin{multicols}{2}
\tableofcontents
\end{multicols}
\vspace{2mm}
\hrule
\end{minipage}

\dancersection{事前課題}{野島 貴英}

今回の事前課題は以下です:
\begin{enumerate}
\item hack time に何をしますか？
\item (任意回答) 本勉強会をどこでお知りになりましたか？
\end{enumerate}
この課題に対して提出いただいた内容は以下です。
\begin{multicols}{2}
{\small
\begin{prework}{ henrich }
  \begin{enumerate}
  \item Q.hack time に何をしますか？\\
    A. バグレポートを読んで、直せそうなところを整理します。
  \end{enumerate}
\end{prework}

\begin{prework}{ wskoka }
  \begin{enumerate}
  \item Q.hack time に何をしますか？\\
    A. tilegxのdebian化
  \end{enumerate}
\end{prework}

\begin{prework}{ rosh }
  \begin{enumerate}
  \item Q.hack time に何をしますか？\\
    A. 自分の担当するパッケージをメインテインする (wide-dhcpv6とadjtimex)
  \end{enumerate}
\end{prework}

\begin{prework}{ tai }
  \begin{enumerate}
  \item Q.hack time に何をしますか？\\
    A. もくもくとパッケージングリハビリとbugs.d.oの落穂ひろいをします。
  \end{enumerate}
\end{prework}

\begin{prework}{ koedoyoshida }
  \begin{enumerate}
  \item Q.hack time に何をしますか？\\
    A. 予定は未定
  \end{enumerate}
\end{prework}

\begin{prework}{ dictoss }
  \begin{enumerate}
  \item Q.hack time に何をしますか？\\
    A. kfreebsd関連作業
  \end{enumerate}
\end{prework}

\begin{prework}{ nametake }
  \begin{enumerate}
  \item Q.hack time に何をしますか？\\
    A. windowsタブレットにdebianインストール
  \end{enumerate}
\end{prework}

\begin{prework}{ yy\_y\_ja\_jp }
  \begin{enumerate}
  \item Q.hack time に何をしますか？\\
    A. DDTSS
  \end{enumerate}
\end{prework}

\begin{prework}{ かじ }
  \begin{enumerate}
  \item Q.hack time に何をしますか？\\
    A. 前回参加できなかったので、今回はwaylandで簡単なサンプルを作ってみたい。\\
参考にしたいのは\\
http://d.hatena.ne.jp/devm33/20140414/1397473785
  \end{enumerate}
\end{prework}

\begin{prework}{ 野島 }
  \begin{enumerate}
  \item Q.hack time に何をしますか？\\
    A. DDTSSとか
  \end{enumerate}
\end{prework}




}
\end{multicols}

\if0
\dancersection{Debian Trivia Quiz}{野島 貴英}

Debianの昨今の話題についてのQuizです。

今回の出題範囲は\url{debian-devel-announce@lists.debian.org} や \url{debian-news@lists.debian.org}に投稿された
内容などからです。

\begin{multicols}{2}
%; whizzy-master ../debianmeetingresume201311.tex
% $B0J>e$N@_Dj$r$7$F$$$k$?$a!"$3$N%U%!%$%k$G(B M-x whizzytex $B$9$k$H!"(Bwhizzytex$B$,MxMQ$G$-$^$9!#(B
%

\santaku
{2015/10/22$B$N(BDPN$B$N%a!<%k$+$iDj4|E*$K$J$,$l$F$$$?$$$/$D$+$N%H%T%C%/$,(BWeb$B$N$_$K7G<($5$l$k$h$&$K$J$j$^$7$?!#$I$3$K7G<($5$l$k!)(B}
{http://www.debian.or.jp/}
{http://www.debian.org/}
{http://bits.debian.org/}
{C}
{2015/10/22$B$N(BDPN$B$N%a!<%k$+$i!"=>MhN.$l$F$$$?(BDPN$B$N%a!<%k$N9=@.$,:~?7$5$l$^$7$?!#$3$A$i$KH<$$!"%;%-%e%j%F%#$K$D$$$F$N%"%J%&%s%9$H!"?7$7$$(BDD/DM$B$N%"%J%&%s%9$O!"(BDPN$B$N%a!<%k$K$O4^$^$l$:!"(BWeb$B%Z!<%8$K7G:\$5$l$k$N$_$H$J$j$^$7$?!#(B}

\santaku
{2015/11/7$B$K$F!"$H$"$k%$%s%9%?%s%H%a%C%;%s%8%c!<MQ%W%m%H%3%k$N%5!<%S%9$,A4(BDebian Developer$B$G;H$($k$h$&$K$J$C$?$H$N%"%J%&%s%9$,$"$j$^$7$?!#%W%m%H%3%k$NL>A0$O<!$N$I$l!)(B}
{XMPP}
{IRC}
{IP Messanger}
{A}
{XMPP$B$O%*!<%W%s$J%$%s%9%?%s%H%a%C%;%s%8%c!<MQ%W%m%H%3%k$N#1$D!#$J$*!"(BDebian$B4X78<T$NMxMQ$K$"$?$C$F>\$7$/$O!"(Bhttp://rtc.debian.org$B$r;2>H!#(B}

\santaku
{2015/10/30$B$K$F!"(BDebian$B$K$F!"(B2$B2sL\$N8xJg$,9T$o$l$?Lr$^$o$j$O<!$N$I$l!)(B}
{2016 DPN}
{technical commitee}
{2016 Debian JP$B2qD9(B}
{B}
{2$BL>$[$I!"(BDebian Developer$B$NJ}$G(Btechnical comittee$B$G3hLv$G$-$kJ}Jg=8$H$N$3$H$G$9!#8=?&(Btechnical comittee$B$N?M$i$,!"(B2015/12/31$B$GG$4|$,@Z$l$F$7$^$&$H$$$&$3$H$G!"$=$l$^$G$K8uJd<T$r5s$2$kI,MW$,$"$k$H$N$3$H$G$9!#(B}



\end{multicols}
\fi

\dancersection{最近のDebian関連のミーティング報告}{野島 貴英}

\subsection{第133回東京エリアDebian勉強会}



% % (query-replace-regexp "<.*?>" "")
% % (query-replace-regexp "^[	 ]\+" "")

%-------------------------------------------------------------------------------
\dancersection{Debian パケージおすすめ}{野島}
%-------------------------------------------------------------------------------

\subsection{はじめに}


%-------------------------------------------------------------------------------
%-------------------------------------------------------------------------------
\dancersection{会場での無線LANのつなぎ方}{野島 貴英,Roger}
%-------------------------------------------------------------------------------
 \subsection{はじめに}

　今回会場側には無線LAN経由のグローバル回線が用意されています。

　以下にDebianマシンでの接続方法を記載します。

 また、自分の環境では違うやり方でつながったという方は、野島まで
教えて下さい。こちらでもノウハウとして溜めていく予定です。

 \subsection{wpasupplicant及び/etc/network/interfacesを利用の場合}

 もっとも良いマニュアルは、/usr/share/doc/wpasupplicant/README.Debian.gz
となります。困った場合はこちらも合わせてご参照下さい。

　以下に/etc/network/interfacesの定義について会場の例を記載します。

\begin{commandline}
$ sudo vi /etc/network/interfaces
-----以下のエントリがなければ追記ここから----------
iface wlan0_debian inet dhcp
     wpa-conf /etc/wpa_supplicant/wpa_supplicant_debian.conf
-----以下のエントリがなければ追記ここまで----------
$ sudo vi /etc/wpa_supplicant/wpa_supplicant_debian.conf
-----以下のエントリを追記ここから----------
network={
     ssid=<<会場のSSID>>
     psk=<<会場のパスワード>>
     scan_ssid=1
}
-----以下のエントリを追記ここまで----------
$ sudo chmod 600 /etc/wpa_supplicant/wpa_supplicant_debian.conf
$ sudo ifup wlan0=wlan0_debian
\end{commandline}
%$

 また、ハマってしまった時のデバッグ方法は、
/usr/share/doc/wpasupplicant/README.Debian.gz中の''4. Trubleshooting''の章が便利です。

 \subsection{その他の無線LAN用パッケージを利用の場合}

　すみません、自分が情報を持たないため、現場で教えて下さい。

\cleartooddpage

\vspace*{15cm}
\hrule
\vspace{2mm}
\includegraphics[width=2cm]{image200502/openlogo-nd.eps}
\noindent \Large \bf Debian 勉強会資料\\
\noindent \normalfont \debmtgyear{}年\debmtgmonth{}月\debmtgdate{}日 \hspace{5mm}  初版第1刷発行\\
\noindent \normalfont 東京エリア Debian 勉強会 （編集・印刷・発行）\\
\hrule

\end{document}
